\scanpage{69}
% archive sheet number: 68

\uline{\unsure{Noter|Notes|Notons}} \((i^*,j^*)\) commutent aux \(\underrightarrow{\lim}\) quelconques \inserted{si elles existent dans \(C',C''\)}, et \(=\) pour que \(\underrightarrow{\lim}_{i}\,(G_i,H_i,\unsure{U_i|u_i|U_i:|\lambda_i})\) existe, il suffit que \(\underrightarrow{\lim}_{i}G_i\) et \(\underrightarrow{\lim}_{i}H_i\) existent, et alors \(\ldots\)

Pour les propriétés d'exactitude \uline{à gauche}, \((i^*,j^*)\) est \deleted{\unsure{exactement|extrêmement}} ``aussi bon que'' \(\varphi:C'\longrightarrow C''\), p.ex. si \(\varphi\) commute aux \(\underleftarrow{\lim}\) finies (resp. quelconques) et si \deleted{\unsure{pour}} \inserted{un} type \deleted{une telle} de \(\underleftarrow{\lim}\) existe \unsure{en|dans} \(C',C''\), alors il existe dans \(C\) et \unsure{les|des} foncteurs \((i^*,j^*)\) y commutent \(\ldots\)
